% ----------------------------------------------------------------------
%
%                            TFMTesis.tex
%
%----------------------------------------------------------------------
%
% Este fichero contiene el "documento maestro" del documento. Lo único
% que hace es configurar el entorno LaTeX e incluir los ficheros .tex
% que contienen cada sección.
%
%----------------------------------------------------------------------
%
% Los ficheros necesarios para este documento son:
%
%       TeXiS/* : ficheros de la plantilla TeXiS.
%       Cascaras/* : ficheros con las partes del documento que no
%          son capítulos ni apéndices (portada, agradecimientos, etc.)
%       Capitulos/*.tex : capítulos de la tesis
%       Apendices/*.tex: apéndices de la tesis
%       constantes.tex: constantes LaTeX
%       config.tex : configuración de la "compilación" del documento
%       guionado.tex : palabras con guiones
%
% Para la bibliografía, además, se necesitan:
%
%       *.bib : ficheros con la información de las referencias
%
% ---------------------------------------------------------------------

\documentclass[12pt,a4paper,twoside]{book}
\usepackage{amsmath}
\usepackage{amsfonts}
\usepackage{amssymb}
\newtheorem{definicion}{Definición}
\newtheorem{proposicion}{Proposición}
\newtheorem{lema}{Lema}
\newtheorem{teorema}{Teorema}
\newtheorem{observacion}{Observación}

%
% Definimos  el   comando  \compilaCapitulo,  que   luego  se  utiliza
% (opcionalmente) en config.tex. Quedaría  mejor si también se definiera
% en  ese fichero,  pero por  el modo  en el  que funciona  eso  no es
% posible. Puedes consultar la documentación de ese fichero para tener
% más  información. Definimos también  \compilaApendice, que  tiene el
% mismo  cometido, pero  que se  utiliza para  compilar  únicamente un
% apéndice.
%
%
% Si  queremos   compilar  solo   una  parte  del   documento  podemos
% especificar mediante  \includeonly{...} qué ficheros  son los únicos
% que queremos  que se incluyan.  Esto  es útil por  ejemplo para sólo
% compilar un capítulo.
%
% El problema es que todos aquellos  ficheros que NO estén en la lista
% NO   se  incluirán...  y   eso  también   afecta  a   ficheros  de
% la plantilla...
%
% Total,  que definimos  una constante  con los  ficheros  que siempre
% vamos a querer compilar  (aquellos relacionados con configuración) y
% luego definimos \compilaCapitulo.
\newcommand{\ficherosBasicosTeXiS}{%
TeXiS/TeXiS_pream,TeXiS/TeXiS_cab,TeXiS/TeXiS_bib,TeXiS/TeXiS_cover%
}
\newcommand{\ficherosBasicosTexto}{%
constantes,guionado,Cascaras/bibliografia,config%
}
\newcommand{\compilaCapitulo}[1]{%
\includeonly{\ficherosBasicosTeXiS,\ficherosBasicosTexto,Capitulos/#1}%
}

\newcommand{\compilaApendice}[1]{%
\includeonly{\ficherosBasicosTeXiS,\ficherosBasicosTexto,Apendices/#1}%
}

%- - - - - - - - - - - - - - - - - - - - - - - - - - - - - - - - - - -
%            Preámbulo del documento. Configuraciones varias
%- - - - - - - - - - - - - - - - - - - - - - - - - - - - - - - - - - -

% Define  el  tipo  de  compilación que  estamos  haciendo.   Contiene
% definiciones  de  constantes que  cambian  el  comportamiento de  la
% compilación. Debe incluirse antes del paquete TeXiS/TeXiS.sty
%---------------------------------------------------------------------
%
%                          config.tex
%
%---------------------------------------------------------------------
%
% Contiene la  definición de constantes  que determinan el modo  en el
% que se compilará el documento.
%
%---------------------------------------------------------------------
%
% En concreto, podemos  indicar si queremos "modo release",  en el que
% no  aparecerán  los  comentarios  (creados  mediante  \com{Texto}  o
% \comp{Texto}) ni los "por  hacer" (creados mediante \todo{Texto}), y
% sí aparecerán los índices. El modo "debug" (o mejor dicho en modo no
% "release" muestra los índices  (construirlos lleva tiempo y son poco
% útiles  salvo  para   la  versión  final),  pero  sí   el  resto  de
% anotaciones.
%
% Si se compila con LaTeX (no  con pdflatex) en modo Debug, también se
% muestran en una esquina de cada página las entradas (en el índice de
% palabras) que referencian  a dicha página (consulta TeXiS_pream.tex,
% en la parte referente a show).
%
% El soporte para  el índice de palabras en  TeXiS es embrionario, por
% lo  que no  asumas que  esto funcionará  correctamente.  Consulta la
% documentación al respecto en TeXiS_pream.tex.
%
%
% También  aquí configuramos  si queremos  o  no que  se incluyan  los
% acrónimos  en el  documento final  en la  versión release.  Para eso
% define (o no) la constante \acronimosEnRelease.
%
% Utilizando \compilaCapitulo{nombre}  podemos también especificar qué
% capítulo(s) queremos que se compilen. Si no se pone nada, se compila
% el documento  completo.  Si se pone, por  ejemplo, 01Introduccion se
% compilará únicamente el fichero Capitulos/01Introduccion.tex
%
% Para compilar varios  capítulos, se separan sus nombres  con comas y
% no se ponen espacios de separación.
%
% En realidad  la macro \compilaCapitulo  está definida en  el fichero
% principal tesis.tex.
%
%---------------------------------------------------------------------


% Comentar la línea si no se compila en modo release.
% TeXiS hará el resto.
% ¡¡¡Si cambias esto, haz un make clean antes de recompilar!!!
\def\release{1}


% Descomentar la linea si se quieren incluir los
% acrónimos en modo release (en modo debug
% no se incluirán nunca).
% ¡¡¡Si cambias esto, haz un make clean antes de recompilar!!!
%\def\acronimosEnRelease{1}


% Descomentar la línea para establecer el capítulo que queremos
% compilar

% \compilaCapitulo{01Introduccion}
% \compilaCapitulo{02EstructuraYGeneracion}
% \compilaCapitulo{03Edicion}
% \compilaCapitulo{04Imagenes}
% \compilaCapitulo{05Bibliografia}
% \compilaCapitulo{06Makefile}

% \compilaApendice{01AsiSeHizo}

% Variable local para emacs, para  que encuentre el fichero maestro de
% compilación y funcionen mejor algunas teclas rápidas de AucTeX
%%%
%%% Local Variables:
%%% mode: latex
%%% TeX-master: "./Tesis.tex"
%%% End:


% Paquete de la plantilla
\usepackage{TeXiS/TeXiS}


% Incluimos el fichero con comandos de constantes
%---------------------------------------------------------------------
%
%                          constantes.tex
%
%---------------------------------------------------------------------
%
% Fichero que  declara nuevos comandos LaTeX  sencillos realizados por
% comodidad en la escritura de determinadas palabras
%
%---------------------------------------------------------------------

%%%%%%%%%%%%%%%%%%%%%%%%%%%%%%%%%%%%%%%%%%%%%%%%%%%%%%%%%%%%%%%%%%%%%%
% Comando: 
%
%       \titulo
%
% Resultado: 
%
% Escribe el título del documento.
%%%%%%%%%%%%%%%%%%%%%%%%%%%%%%%%%%%%%%%%%%%%%%%%%%%%%%%%%%%%%%%%%%%%%%
\def\titulo{\textsc{TeXiS}: Una plantilla de \LaTeX\
  para Tesis y otros documentos}

%%%%%%%%%%%%%%%%%%%%%%%%%%%%%%%%%%%%%%%%%%%%%%%%%%%%%%%%%%%%%%%%%%%%%%
% Comando: 
%
%       \autor
%
% Resultado: 
%
% Escribe el autor del documento.
%%%%%%%%%%%%%%%%%%%%%%%%%%%%%%%%%%%%%%%%%%%%%%%%%%%%%%%%%%%%%%%%%%%%%%
\def\autor{Marco Antonio y Pedro Pablo G\'omez Mart\'in}

% Variable local para emacs, para  que encuentre el fichero maestro de
% compilación y funcionen mejor algunas teclas rápidas de AucTeX

%%%
%%% Local Variables:
%%% mode: latex
%%% TeX-master: "tesis.tex"
%%% End:


% Sacamos en el log de la compilación el copyright
%\typeout{Copyright Marco Antonio and Pedro Pablo Gomez Martin}

%
% "Metadatos" para el PDF
%
\ifpdf\hypersetup{%
    pdftitle = {\titulo},
    pdfsubject = {Plantilla de Tesis},
    pdfkeywords = {Plantilla, LaTeX, tesis, trabajo de
      investigación, trabajo de Master},
    pdfauthor = {\textcopyright\ \autor},
    pdfcreator = {\LaTeX\ con el paquete \flqq hyperref\frqq},
    pdfproducer = {pdfeTeX-0.\the\pdftexversion\pdftexrevision},
    }
    \pdfinfo{/CreationDate (\today)}
\fi


%- - - - - - - - - - - - - - - - - - - - - - - - - - - - - - - - - - -
%                        Documento
%- - - - - - - - - - - - - - - - - - - - - - - - - - - - - - - - - - -
\begin{document}

% Incluimos el  fichero de definición de guionado  de algunas palabras
% que LaTeX no ha dividido como debería
%----------------------------------------------------------------
%
%                          guionado.tex
%
%----------------------------------------------------------------
%
% Fichero con algunas divisiones de palabras que LaTeX no
% hace correctamente si no se le da alguna ayuda.
%
%----------------------------------------------------------------

\hyphenation{
% a
abs-trac-to
abs-trac-tos
abs-trac-ta
abs-trac-tas
ac-tua-do-res
a-gra-de-ci-mien-tos
ana-li-za-dor
an-te-rio-res
an-te-rior-men-te
apa-rien-cia
a-pro-pia-do
a-pro-pia-dos
a-pro-pia-da
a-pro-pia-das
a-pro-ve-cha-mien-to
a-que-llo
a-que-llos
a-que-lla
a-que-llas
a-sig-na-tu-ra
a-sig-na-tu-ras
a-so-cia-da
a-so-cia-das
a-so-cia-do
a-so-cia-dos
au-to-ma-ti-za-do
% b
batch
bi-blio-gra-fía
bi-blio-grá-fi-cas
bien
bo-rra-dor
boo-l-ean-expr
% c
ca-be-ce-ra
call-me-thod-ins-truc-tion
cas-te-lla-no
cir-cuns-tan-cia
cir-cuns-tan-cias
co-he-ren-te
co-he-ren-tes
co-he-ren-cia
co-li-bri
co-men-ta-rio
co-mer-cia-les
co-no-ci-mien-to
cons-cien-te
con-si-de-ra-ba
con-si-de-ra-mos
con-si-de-rar-se
cons-tan-te
cons-trucción
cons-tru-ye
cons-tru-ir-se
con-tro-le
co-rrec-ta-men-te
co-rres-pon-den
co-rres-pon-dien-te
co-rres-pon-dien-tes
co-ti-dia-na
co-ti-dia-no
crean
cris-ta-li-zan
cu-rri-cu-la
cu-rri-cu-lum
cu-rri-cu-lar
cu-rri-cu-la-res
% d
de-di-ca-do
de-di-ca-dos
de-di-ca-da
de-di-ca-das
de-rro-te-ro
de-rro-te-ros
de-sa-rro-llo
de-sa-rro-llos
de-sa-rro-lla-do
de-sa-rro-lla-dos
de-sa-rro-lla-da
de-sa-rro-lla-das
de-sa-rro-lla-dor
de-sa-rro-llar
des-cri-bi-re-mos
des-crip-ción
des-crip-cio-nes
des-cri-to
des-pués
de-ta-lla-do
de-ta-lla-dos
de-ta-lla-da
de-ta-lla-das
di-a-gra-ma
di-a-gra-mas
di-se-ños
dis-po-ner
dis-po-ni-bi-li-dad
do-cu-men-ta-da
do-cu-men-to
do-cu-men-tos
% e
edi-ta-do
e-du-ca-ti-vo
e-du-ca-ti-vos
e-du-ca-ti-va
e-du-ca-ti-vas
e-la-bo-ra-do
e-la-bo-ra-dos
e-la-bo-ra-da
e-la-bo-ra-das
es-co-llo
es-co-llos
es-tu-dia-do
es-tu-dia-dos
es-tu-dia-da
es-tu-dia-das
es-tu-dian-te
e-va-lua-cio-nes
e-va-lua-do-res
exis-ten-tes
exhaus-ti-va
ex-pe-rien-cia
ex-pe-rien-cias
% f
for-ma-li-za-do
% g
ge-ne-ra-ción
ge-ne-ra-dor
ge-ne-ra-do-res
ge-ne-ran
% h
he-rra-mien-ta
he-rra-mien-tas
% i
i-dio-ma
i-dio-mas
im-pres-cin-di-ble
im-pres-cin-di-bles
in-de-xa-do
in-de-xa-dos
in-de-xa-da
in-de-xa-das
in-di-vi-dual
in-fe-ren-cia
in-fe-ren-cias
in-for-ma-ti-ca
in-gre-dien-te
in-gre-dien-tes
in-me-dia-ta-men-te
ins-ta-la-do
ins-tan-cias
% j
% k
% l
len-gua-je
li-be-ra-to-rio
li-be-ra-to-rios
li-be-ra-to-ria
li-be-ra-to-rias
li-mi-ta-do
li-te-ra-rio
li-te-ra-rios
li-te-ra-ria
li-te-ra-rias
lo-tes
% m
ma-ne-ra
ma-nual
mas-que-ra-de
ma-yor
me-mo-ria
mi-nis-te-rio
mi-nis-te-rios
mo-de-lo
mo-de-los
mo-de-la-do
mo-du-la-ri-dad
mo-vi-mien-to
% n
na-tu-ral
ni-vel
nues-tro
% o
obs-tan-te
o-rien-ta-do
o-rien-ta-dos
o-rien-ta-da
o-rien-ta-das
% p
pa-ra-le-lo
pa-ra-le-la
par-ti-cu-lar
par-ti-cu-lar-men-te
pe-da-gó-gi-ca
pe-da-gó-gi-cas
pe-da-gó-gi-co
pe-da-gó-gi-cos
pe-rio-di-ci-dad
per-so-na-je
plan-te-a-mien-to
plan-te-a-mien-tos
po-si-ción
pre-fe-ren-cia
pre-fe-ren-cias
pres-cin-di-ble
pres-cin-di-bles
pri-me-ra
pro-ble-ma
pro-ble-mas
pró-xi-mo
pu-bli-ca-cio-nes
pu-bli-ca-do
% q
% r
rá-pi-da
rá-pi-do
ra-zo-na-mien-to
ra-zo-na-mien-tos
re-a-li-zan-do
re-fe-ren-cia
re-fe-ren-cias
re-fe-ren-cia-da
re-fe-ren-cian
re-le-van-tes
re-pre-sen-ta-do
re-pre-sen-ta-dos
re-pre-sen-ta-da
re-pre-sen-ta-das
re-pre-sen-tar-lo
re-qui-si-to
re-qui-si-tos
res-pon-der
res-pon-sa-ble
% s
se-pa-ra-do
si-guien-do
si-guien-te
si-guien-tes
si-guie-ron
si-mi-lar
si-mi-la-res
si-tua-ción
% t
tem-pe-ra-ments
te-ner
trans-fe-ren-cia
trans-fe-ren-cias
% u
u-sua-rio
Unreal-Ed
% v
va-lor
va-lo-res
va-rian-te
ver-da-de-ro
ver-da-de-ros
ver-da-de-ra
ver-da-de-ras
ver-da-de-ra-men-te
ve-ri-fi-ca
% w
% x
% y
% z
}
% Variable local para emacs, para que encuentre el fichero
% maestro de compilación
%%%
%%% Local Variables:
%%% mode: latex
%%% TeX-master: "./Tesis.tex"
%%% End:


% Marcamos  el inicio  del  documento para  la  numeración de  páginas
% (usando números romanos para esta primera fase).
\frontmatter
\pagestyle{empty}

%---------------------------------------------------------------------
%
%                          configCover.tex
%
%---------------------------------------------------------------------
%
% cover.tex
% Copyright 2009 Marco Antonio Gomez-Martin, Pedro Pablo Gomez-Martin
%
% This file belongs to the TeXiS manual, a LaTeX template for writting
% Thesis and other documents. The complete last TeXiS package can
% be obtained from http://gaia.fdi.ucm.es/projects/texis/
%
% Although the TeXiS template itself is distributed under the 
% conditions of the LaTeX Project Public License
% (http://www.latex-project.org/lppl.txt), the manual content
% uses the CC-BY-SA license that stays that you are free:
%
%    - to share & to copy, distribute and transmit the work
%    - to remix and to adapt the work
%
% under the following conditions:
%
%    - Attribution: you must attribute the work in the manner
%      specified by the author or licensor (but not in any way that
%      suggests that they endorse you or your use of the work).
%    - Share Alike: if you alter, transform, or build upon this
%      work, you may distribute the resulting work only under the
%      same, similar or a compatible license.
%
% The complete license is available in
% http://creativecommons.org/licenses/by-sa/3.0/legalcode
%
%---------------------------------------------------------------------
%
% Fichero que contiene la configuración de la portada y de la 
% primera hoja del documento.
%
%---------------------------------------------------------------------


% Pueden configurarse todos los elementos del contenido de la portada
% utilizando comandos.

%%%%%%%%%%%%%%%%%%%%%%%%%%%%%%%%%%%%%%%%%%%%%%%%%%%%%%%%%%%%%%%%%%%%%%
% Título del documento:
% \tituloPortada{titulo}
% Nota:
% Si no se define se utiliza el del \titulo. Este comando permite
% cambiar el título de forma que se especifiquen dónde se quieren
% los retornos de carro cuando se utilizan fuentes grandes.
%%%%%%%%%%%%%%%%%%%%%%%%%%%%%%%%%%%%%%%%%%%%%%%%%%%%%%%%%%%%%%%%%%%%%%
\tituloPortada{%
Estudio del problema del plegado de proteínas y otros problemas de biología
}


%%%%%%%%%%%%%%%%%%%%%%%%%%%%%%%%%%%%%%%%%%%%%%%%%%%%%%%%%%%%%%%%%%%%%%
% Título del documento en inglés:
% \tituloPortadaEng{titulo}
% Nota:
% Si no se define se utiliza el del \titulo. Este comando permite
% cambiar el título de forma que se especifiquen dónde se quieren
% los retornos de carro cuando se utilizan fuentes grandes.
%%%%%%%%%%%%%%%%%%%%%%%%%%%%%%%%%%%%%%%%%%%%%%%%%%%%%%%%%%%%%%%%%%%%%%
\tituloPortadaEng{%
Analysis of the protein folding problem and other problems from biology
}

%%%%%%%%%%%%%%%%%%%%%%%%%%%%%%%%%%%%%%%%%%%%%%%%%%%%%%%%%%%%%%%%%%%%%%
% Autor del documento:
% \autorPortada{Nombre}
% Se utiliza en la portada y en el valor por defecto del
% primer subtítulo de la segunda portada.
%%%%%%%%%%%%%%%%%%%%%%%%%%%%%%%%%%%%%%%%%%%%%%%%%%%%%%%%%%%%%%%%%%%%%%
\autorPortada{Javier Jiménez Arenas}

%%%%%%%%%%%%%%%%%%%%%%%%%%%%%%%%%%%%%%%%%%%%%%%%%%%%%%%%%%%%%%%%%%%%%%
% Fecha de publicación:
% \fechaPublicacion{Fecha}
% Puede ser vacío. Aparece en la última línea de ambas portadas
%%%%%%%%%%%%%%%%%%%%%%%%%%%%%%%%%%%%%%%%%%%%%%%%%%%%%%%%%%%%%%%%%%%%%%
% Descomentar para que ponga siempre la fecha actual
%\fechaPublicacion{\today}
\fechaPublicacion{\textcolor{red}{DIA de MES de AÑO}}

%%%%%%%%%%%%%%%%%%%%%%%%%%%%%%%%%%%%%%%%%%%%%%%%%%%%%%%%%%%%%%%%%%%%%%
% Imagen de la portada (y escala)
% \imagenPortada{Fichero}
% \escalaImagenPortada{Numero}
% Si no se especifica, se utiliza la imagen TODO.pdf
%%%%%%%%%%%%%%%%%%%%%%%%%%%%%%%%%%%%%%%%%%%%%%%%%%%%%%%%%%%%%%%%%%%%%%
% imagen en blanco y negro
%\imagenPortada{Imagenes/Vectorial/escudoUCM}
%imagen en color
\imagenPortada{Imagenes/Bitmap/escudoUCMcolor}
\escalaImagenPortada{.2}

%%%%%%%%%%%%%%%%%%%%%%%%%%%%%%%%%%%%%%%%%%%%%%%%%%%%%%%%%%%%%%%%%%%%%%
% Tipo de documento.
% \tipoDocumento{Tipo}
% Para el texto justo debajo del escudo.
% Si no se indica, se utiliza "TESIS DOCTORAL".
%%%%%%%%%%%%%%%%%%%%%%%%%%%%%%%%%%%%%%%%%%%%%%%%%%%%%%%%%%%%%%%%%%%%%%
\tipoDocumento{Trabajo de Fin de Grado}

%%%%%%%%%%%%%%%%%%%%%%%%%%%%%%%%%%%%%%%%%%%%%%%%%%%%%%%%%%%%%%%%%%%%%%
% Institución/departamento asociado al documento.
% \institucion{Nombre}
% Puede tener varias líneas. Se utiliza en las dos portadas.
% Si no se indica aparecerá vacío.
%%%%%%%%%%%%%%%%%%%%%%%%%%%%%%%%%%%%%%%%%%%%%%%%%%%%%%%%%%%%%%%%%%%%%%
\institucion{%
Doble Grado en Ingeniería Informática y Matemáticas\\[0.2em]
Facultad de Informática\\[0.2em]
Universidad Complutense de Madrid
}

%%%%%%%%%%%%%%%%%%%%%%%%%%%%%%%%%%%%%%%%%%%%%%%%%%%%%%%%%%%%%%%%%%%%%%
% Director del trabajo.
% \directorPortada{Nombre}
% Se utiliza para el valor por defecto del segundo subtítulo, donde
% se indica quién es el director del trabajo.
% Si se fuerza un subtítulo distinto, no hace falta definirlo.
%%%%%%%%%%%%%%%%%%%%%%%%%%%%%%%%%%%%%%%%%%%%%%%%%%%%%%%%%%%%%%%%%%%%%%
\directorPortada{Ismael Rodríguez Laguna\\Fernando Rubio Diez}


%%%%%%%%%%%%%%%%%%%%%%%%%%%%%%%%%%%%%%%%%%%%%%%%%%%%%%%%%%%%%%%%%%%%%%
% Colaborador en la dirección del trabajo.
% \colaboradorPortada{Nombre}
% Se utiliza para el valor por defecto del segundo subtítulo, donde
% se indica quién es el colaborador en la dirección del trabajo.
% Si se fuerza un subtítulo distinto, no hace falta definirlo.
%%%%%%%%%%%%%%%%%%%%%%%%%%%%%%%%%%%%%%%%%%%%%%%%%%%%%%%%%%%%%%%%%%%%%%
\colaboradorPortada{\textcolor{red}{Colaborador 1\\Colaborador 2}}


%%%%%%%%%%%%%%%%%%%%%%%%%%%%%%%%%%%%%%%%%%%%%%%%%%%%%%%%%%%%%%%%%%%%%%
% Texto del primer subtítulo de la segunda portada.
% \textoPrimerSubtituloPortada{Texto}
% Para configurar el primer "texto libre" de la segunda portada.
% Si no se especifica se indica "Memoria que presenta para optar al
% título de Doctor en Informática" seguido del \autorPortada.
%%%%%%%%%%%%%%%%%%%%%%%%%%%%%%%%%%%%%%%%%%%%%%%%%%%%%%%%%%%%%%%%%%%%%%
\textoPrimerSubtituloPortada{%
\textbf{Trabajo de Fin de Grado en \textcolor{red}{Ingeniería Informática}}\\ [0.3em]
}

%%%%%%%%%%%%%%%%%%%%%%%%%%%%%%%%%%%%%%%%%%%%%%%%%%%%%%%%%%%%%%%%%%%%%%
% Texto del segundo subtítulo de la segunda portada.
% \textoSegundoSubtituloPortada{Texto}
% Para configurar el segundo "texto libre" de la segunda portada.
% Si no se especifica se indica "Dirigida por el Doctor" seguido
% del \directorPortada.
%%%%%%%%%%%%%%%%%%%%%%%%%%%%%%%%%%%%%%%%%%%%%%%%%%%%%%%%%%%%%%%%%%%%%%
\textoSegundoSubtituloPortada{%
\textbf{Convocatoria: }\textit{\textcolor{red}{Febrero/Junio/Septiembre} \the\year}%\\[0.2em]
%\textbf{Calificación: }\textit{\textcolor{red}{Nota}}
}

%%%%%%%%%%%%%%%%%%%%%%%%%%%%%%%%%%%%%%%%%%%%%%%%%%%%%%%%%%%%%%%%%%%%%%
% \explicacionDobleCara
% Si se utiliza, se aclara que el documento está preparado para la
% impresión a doble cara.
%%%%%%%%%%%%%%%%%%%%%%%%%%%%%%%%%%%%%%%%%%%%%%%%%%%%%%%%%%%%%%%%%%%%%%
%\explicacionDobleCara

%%%%%%%%%%%%%%%%%%%%%%%%%%%%%%%%%%%%%%%%%%%%%%%%%%%%%%%%%%%%%%%%%%%%%%
% \isbn
% Si se utiliza, aparecerá el ISBN detrás de la segunda portada.
%%%%%%%%%%%%%%%%%%%%%%%%%%%%%%%%%%%%%%%%%%%%%%%%%%%%%%%%%%%%%%%%%%%%%%
%\isbn{978-84-692-7109-4}


%%%%%%%%%%%%%%%%%%%%%%%%%%%%%%%%%%%%%%%%%%%%%%%%%%%%%%%%%%%%%%%%%%%%%%
% \copyrightInfo
% Si se utiliza, aparecerá información de los derechos de copyright
% detrás de la segunda portada.
%%%%%%%%%%%%%%%%%%%%%%%%%%%%%%%%%%%%%%%%%%%%%%%%%%%%%%%%%%%%%%%%%%%%%%
%\copyrightInfo{\autor}


%%
%% Creamos las portadas
%%
\makeCover

% Variable local para emacs, para que encuentre el fichero
% maestro de compilación
%%%
%%% Local Variables:
%%% mode: latex
%%% TeX-master: "../Tesis.tex"
%%% End:

%\include{Cascaras/autorizacion}
% +--------------------------------------------------------------------+
% | Dedication Page (Optional)
% +--------------------------------------------------------------------+

\chapter*{Dedicatoria}

\begin{flushright}
\begin{minipage}[c]{8.5cm}
\flushright{\textit{Aun es pronto para esto}}
\end{minipage}
\end{flushright}
% +--------------------------------------------------------------------+
% | Acknowledgements Page (Optional)                                   |
% +--------------------------------------------------------------------+

\chapter*{Agradecimientos}

Aun es pronto para esto












\chapter*{Resumen}

\section*{\tituloPortadaVal}

Un resumen en castellano de media página, incluyendo el título en castellano. A continuación, se escribirá una lista de no más de 10 palabras clave.


\section*{Palabras clave}
   
\noindent Máximo 10 palabras clave separadas por comas

   



\begin{otherlanguage}{english}
\chapter*{Abstract}

\section*{\tituloPortadaEngVal}

An abstract in English, half a page long, including the title in English. Below, a list with no more than 10 keywords.


\section*{Keywords}

\noindent 10 keywords max., separated by commas.




% Si el trabajo se escribe en inglés, comentar esta línea y descomentar
% otra igual que hay justo antes de \end{document}
\end{otherlanguage}

\ifx\generatoc\undefined
\else
%---------------------------------------------------------------------
%
%                          TeXiS_toc.tex
%
%---------------------------------------------------------------------
%
% TeXiS_toc.tex
% Copyright 2009 Marco Antonio Gomez-Martin, Pedro Pablo Gomez-Martin
%
% This file belongs to TeXiS, a LaTeX template for writting
% Thesis and other documents. The complete last TeXiS package can
% be obtained from http://gaia.fdi.ucm.es/projects/texis/
%
% This work may be distributed and/or modified under the
% conditions of the LaTeX Project Public License, either version 1.3
% of this license or (at your option) any later version.
% The latest version of this license is in
%   http://www.latex-project.org/lppl.txt
% and version 1.3 or later is part of all distributions of LaTeX
% version 2005/12/01 or later.
%
% This work has the LPPL maintenance status `maintained'.
% 
% The Current Maintainers of this work are Marco Antonio Gomez-Martin
% and Pedro Pablo Gomez-Martin
%
%---------------------------------------------------------------------
%
% Contiene  los  comandos  para  generar los  índices  del  documento,
% entendiendo por índices las tablas de contenidos.
%
% Genera  el  índice normal  ("tabla  de  contenidos"),  el índice  de
% figuras y el de tablas. También  crea "marcadores" en el caso de que
% se esté compilando con pdflatex para que aparezcan en el PDF.
%
%---------------------------------------------------------------------


% Primero un poquito de configuración...


% Pedimos que inserte todos los epígrafes hasta el nivel \subsection en
% la tabla de contenidos.
\setcounter{tocdepth}{2} 

% Le  pedimos  que nos  numere  todos  los  epígrafes hasta  el  nivel
% \subsubsection en el cuerpo del documento.
\setcounter{secnumdepth}{3} 


% Creamos los diferentes índices.

% Lo primero un  poco de trabajo en los marcadores  del PDF. No quiero
% que  salga una  entrada  por cada  índice  a nivel  0...  si no  que
% aparezca un marcador "Índices", que  tenga dentro los otros tipos de
% índices.  Total, que creamos el marcador "Índices".
% Antes de  la creación  de los índices,  se añaden los  marcadores de
% nivel 1.

\ifpdf
   \pdfbookmark{Índices}{indices}
\fi

% Tabla de contenidos.
%
% La  inclusión  de '\tableofcontents'  significa  que  en la  primera
% pasada  de  LaTeX  se  crea   un  fichero  con  extensión  .toc  con
% información sobre la tabla de contenidos (es conceptualmente similar
% al  .bbl de  BibTeX, creo).  En la  segunda ejecución  de  LaTeX ese
% documento se utiliza para  generar la verdadera página de contenidos
% usando la  información sobre los  capítulos y demás guardadas  en el
% .toc
\ifpdf
   \pdfbookmark[1]{Tabla de Contenidos}{tabla de contenidos}
\fi

\cabeceraEspecial{\'Indice}

\tableofcontents

\newpage 

% Índice de figuras
%
% La idea es semejante que para  el .toc del índice, pero ahora se usa
% extensión .lof (List Of Figures) con la información de las figuras.

\ifpdf
   \pdfbookmark[1]{Índice de figuras}{indice de figuras}
\fi

\cabeceraEspecial{\'Indice de figuras}

\listoffigures

\newpage

% Índice de tablas
% Como antes, pero ahora .lot (List Of Tables)

\ifpdf
   \pdfbookmark[1]{Índice de tablas}{indice de tablas}
\fi

\cabeceraEspecial{\'Indice de tablas}

\listoftables

\newpage

% Variable local para emacs, para  que encuentre el fichero maestro de
% compilación y funcionen mejor algunas teclas rápidas de AucTeX

%%%
%%% Local Variables:
%%% mode: latex
%%% TeX-master: "../Tesis.tex"
%%% End:

\fi

% Marcamos el  comienzo de  los capítulos (para  la numeración  de las
% páginas) y ponemos la cabecera normal
\mainmatter

\pagestyle{fancy}
\restauraCabecera

\setlength{\parskip}{3mm}

\chapter{Introducción}
\label{cap:introduccion}

%\chapterquote{Frase célebre dicha por alguien inteligente}{Autor}
\section{Antecedentes}
La bioinformática es una disciplina científica que persigue la investigación, desarrollo y aplicación de métodos computacionales y aproximaciones en problemas de índole biológico o relacionados con la salud.

Entre estos problemas, uno de los más destacados es el problema de predicción del plegado de proteínas. Este problema consiste en predecir, dada una cadena de aminoácidos, la estructura tridimensional que tendrá la proteína resultante. La importancia de este problema radica en el hecho de que esta estructura es de vital importancia para el correcto desarrollo de las funciones asociadas a la proteína. Tanto es así que se ha probado que un plegamiento incorrecto de proteínas puede provocar enfermedades. Además, poder predecir la estructura ayudaría a avanzar en el desarrollo de proteínas artificiales en laboratorios en el ámbito de la salud.

En el marco más teórico, se han desarrollado una variedad de modelos matemáticos con los que estudiar la complejidad computacional del problema, esto es, la viabilidad de encontrar la solucion al problema en tiempo razonable. Algunos resultados apuntarían a que este problema resultaría ser $\mathcal{NP}$-completo. Sin embargo, hay discusiones sobre el correcto modelizado del problema que podría arrojar una falsa complejidad sobre el mismo. De cualquier manera, con otros modelos más completos se intuye la $\mathcal{NP}$-completitud del problema. Esto también nos lleva a interesarnos por la aproximabilidad del problema, es decir, la viabilidad de obtener soluciones cercanas a la óptima en un tiempo razonable. Más adelante daremos una definición más específica sobre los conceptos de complejidad y aproximabilidad mencionados anteriormente.

\textcolor{red}{Aún habrá que añadir cosas.}

\newpage

\section{Objetivos del trabajo}
\textcolor{red}{Aún no tengo esto nada claro. Habrá que cambiarlo más adelante. Digamos el estudio de varios modelos y la definición de alguno propio y hablar sobre la aproximabilidad de los distintos modelos. No tanto sobre la complejidad.}

\section{Plan de trabajo}
\textcolor{red}{De nuevo, muy pronto para escribirlo.}

\chapter{Complejidad y Aproximabilidad}
En este capítulo haremos una introducción a las Teorías de Complejidad y Aproximabilidad. Definiremos los conceptos más importantes que serán la base de este trabajo incluyendo los tipos de problemas que se estudian y las clasificaciones de los mismos.

En este sentido, cubriremos los problemas de decisión y las clases $\mathcal{P}$ y $\mathcal{NP}$, así como los conceptos de $\mathcal{NP}$-completitud y $\mathcal{NP}$-dureza, pasando por las reducciones de problemas que nos sirven para clasificar. Por otro lado, hablaremos de los problemas de optimización y de las clases de problemas $\mathcal{PO}$, $\mathcal{NPO}$ y muchas otras clases intermedias.

\section{Problemas de Decisión y Clases de Complejidad}
Los problemas de decisión podrían verse intuitivamente como aquellos que dada una entrada devuelven una respuesta de sí o no. Este tipo de problemas son la base que usamos para hablar de computabilidad y complejidad. Definimos a continuación los conceptos básicos a tener en cuenta para hablar de complejidad, que se basan en los definidos en \cite{compcomplexbarak}.

\begin{definicion}
Dada una función booleana $f: \{0,1\}^* \rightarrow \{0,1\}$, llamamos \textbf{problema de decisión} al conjunto $L_f = \{x \in \{0,1\}^{*} \ | \ f(x) = 1\}$.
\end{definicion}

\begin{definicion}
Suponemos que tenemos cintas de longitud infinita donde podemos leer y escribir símbolos moviéndonos solo a la derecha o la izquierda y un registro donde guardar el último estado de la máquina. Definimos una \textbf{Máquina de Turing} $M$ como la tupla $(\Gamma, Q, \delta)$, donde:
\begin{itemize}
    \item[$\bullet$] $\Gamma$ es un conjunto de símbolos que pueden ser escritos en las cintas de $M$. Para nuestro caso podemos asumir que $\Gamma$ contiene un símbolo vacío, un símbolo de inicio y los números 0 y 1.
    \item[$\bullet$] $Q$ es un conjunto de posibles estados que pueden guardarse en el registro de $M$. Asumimos que el registro incialmente contiene el estado inicial y que existe un estado final de parada.
    \item[$\bullet$] $\delta$ es una función $\delta: Q \times \Gamma^{k} \rightarrow Q \times \Gamma^{k-1} \times \{L,S,R\}^{k}$ que describe las reglas que $M$ seguirá en cada paso y en cada cinta (L significa left, R significa right y S significa stay. Corresponde en español a moverse a la izquierda, la derecha o no moverse). Esta función se llama función de transición de $M$.
\end{itemize}
\end{definicion}
Existen resultados importantes sobre máquinas de Turing que no vamos a probar en este trabajo pero que son necesarios para definir conceptos clave del trabajo. Estos dos resultados son:

\begin{proposicion}
Las siguientes afirmaciones son ciertas:
\begin{itemize}
    \item[$\bullet$] Cada cadena $x \in \{0,1\}^*$ representa alguna máquina de Turing.
    \item[$\bullet$] Cada máquina de Turing puede ser representada con infinitas cadenas distintas de $\{0,1\}^*$
\end{itemize}
\end{proposicion}

\begin{proposicion}
Existe una \textbf{Máquina de Turing Universal} $\mathcal{U}$ capaz de simular el comportamiento de cualquier otra dada su definición binaria. Esto lo representamos como $\mathcal{U}(x,\alpha) = M_{\alpha}(x)$.
\end{proposicion}

\begin{definicion}[Computar una función y tiempo de cómputo]
Dadas dos funciones $f: \{0,1\}^* \rightarrow \{0,1\}^*$ y $T: \mathbb{N} \rightarrow \mathbb{N}$ y una máquina de Turing $M$, decimos que $M$ computa $f$ en tiempo $T(n)$ si $\forall x \in \{0,1\}^*$ y suponiendo que $M$ está en su configuración inicial con $x$ en la cinta de entrada entonces la máquina llega al estado de parada (a partir de ahora diremos que la máquina termina) después de a lo sumo $T(|x|)$ pasos con $f(x)$ escrito en la cinta de salida. \\
Decimos que $M$ computa $f$ si computa $f$ en tiempo $T(n)$ para alguna función $T: \mathbb{N} \rightarrow \mathbb{N}$.
\end{definicion}

Es importante ver que estas definiciones tratan los casos peores, pues el tiempo de computabilidad es una cota superior para la instancia que más pasos necesite para terminar. Con estas definiciones en mente ya podemos hablar de las clases de complejidad que dan lugar al famoso problema de P vs NP, uno de los llamados problemas del milenio.

\begin{definicion}[Clase $\mathcal{P}$]
La \textbf{clase $\mathcal{P}$} es el conjunto de todos los problemas de decisión tales que para cada $f$ problema de decisión existe $p$ polinomio por el cual $f$ es computable en tiempo $p$.
\end{definicion}

Esta clase corresponde a la de los problemas ``resolubles'' en tiempo polinómico. Ahora hablaremos de la clase $\mathcal{NP}$, que es intuitivamente la de los problemas con solución ``verificable'' en tiempo polinómico. 
\begin{definicion}[Clase $\mathcal{NP}$]
Un problema de decisión $L_f$ está en $\mathcal{NP}$ si existe un polinomio $p: \mathbb{N} \rightarrow \mathbb{N}$ y una máquina de Turing M de tiempo polinómico tal que
$$\forall x \in \{0,1\}^{*}, \ x \in L_f \Leftrightarrow \exists u \in \{0,1\}^{p(|x|)} \ tal \ que \ M(x, u) = 1$$
\end{definicion}

Cabe destacar que conseguir una solución en tiempo polinómico implica que la verificación también se hará en tiempo polinómico, por lo que se tiene que $\mathcal{P} \subseteq \mathcal{NP}$ de forma sencilla. Sin embargo, que la inclusión sea estricta o no es objeto de fuerte debate, pues las consecuencias de que ambos conjuntos fuesen iguales serían trascendentales.

\begin{definicion}
    Decimos que un problema de decisión $A$ es reducible Karp en tiempo polinomial a otro problema de decisión B, y denotamos $A \leq_{p} B$, si existe una función $f: \{0,1\}^* \rightarrow \{0,1\}^*$ computable en tiempo polinómico tal que $$\forall x \in \{0,1\}^*, \ x \in A \Leftrightarrow f(x) \in B$$
\end{definicion}
\begin{definicion}[$\mathcal{NP}$-dureza]
Decimos que un problema B es $\mathcal{NP}$-duro si $A \leq_{p} B \ \forall A \in \mathcal{NP}$.
\end{definicion}

\begin{definicion}[$\mathcal{NP}$-completitud]
    Decimos que un problema B es $\mathcal{NP}$-completo si es $\mathcal{NP}$-duro y $B \in \mathcal{NP}$.
\end{definicion}

\section{Problemas de Optimización y Clases de Aproximabilidad}
Paralelamente a los problemas de decisión podemos definir los problemas de optimización, en los cuales la idea es obtener el valor máximo o mínimo posible para una función. De aquí podemos obtener clases de problemas en función de lo factible que es encontrar soluciones cercanas a la óptima. Las definiciones que se especifican a continuación se basan en los artículos \cite{aproxpaschos} y \cite{aproxcres}.
\begin{definicion}
Un \textbf{problema de optimización} es una 4-tupla P=($\mathcal{I}$, Sol, m, goal) donde:
\begin{itemize}
    \item[$\bullet$] $\mathcal{I}$ es el conjunto de instancias del problema.
    \item[$\bullet$] Dada una instancia $I \in \mathcal{I}$, $Sol(I)$ es el conjunto de soluciones factibles para $I$.
    \item[$\bullet$] $m: \mathcal{I} \times Sol \rightarrow \mathbb{N}$ es la función de valor objetivo, que asocia a cada instancia y cada solución un valor numérico, que podemos suponer natural. En este caso por abuso de notación denotamos $Sol = \displaystyle{\bigcup_{I \in \mathcal{I}} Sol(I)}$.
    \item[$\bullet$] $goal \in \{max, min\}$ indica si el problema es de maximización o de minimización. 
\end{itemize}
\end{definicion}

\begin{definicion}
    Decimos que un problema de optimización pertenece a la \textbf{clase $\mathcal{NPO}$} si:
    \begin{itemize}
        \item[$\bullet$] El conjunto de instancias $\mathcal{I}$ es reconocible en tiempo polinómico. Es decir, dado $I$, el problema de decisión que resuelve si $I \in \mathcal{I}$ pertenece a la clase $\mathcal{P}$.
        \item[$\bullet$] Las soluciones factibles de una instancia de un problema son cortas. Es decir, dado $I \in \mathcal{I}$, $\exists \ p$ un polinomio tal que $\forall y \in Sol(I)$ se tiene $y \leq p(I)$. Además, se puede calcular la factibilidad de una solución corta en tiempo polinómico. Esto es, dado un $y \leq p(I)$, es decidible en tiempo polinómico si $y \in Sol(I)$.
        \item[$\bullet$] $\forall I \in \mathcal{I}, \ \forall S \in Sol(I), \ m(I,S)$ se puede calcular en tiempo polinómico.
    \end{itemize}
\end{definicion}

\begin{definicion}
    Decimos que un problema de optimización pertenece a la clase $\mathcal{PO}$ si:
    \begin{itemize}
        \item[$\bullet$] Pertenece a la clase $\mathcal{NPO}$.
        \item[$\bullet$] $\forall I \in \mathcal{I}$ se puede encontrar una solución óptima al problema en tiempo polinómico.
    \end{itemize}
\end{definicion}

\begin{definicion}
    Llamaremos \textbf{algoritmo de aproximación} A para un problema de optimización P es una función que dada una instancia $I \in \mathcal{I}$ devuelve una solución factible $A(I) = S \in Sol$. 
\end{definicion}
Para los problemas de optimización, se habla de su aproximabilidad en función de un ratio o cociente. Existen dos paradigmas principales a seguir. Para estas definiciones, tomaremos $\omega(I)$ como el valor de la peor solución para la instancia $I$, $m_{A}(I, S)$ como el valor de la solución aproximada $A(I) = S$, que será la solución obtenida por un algoritmo de aproximación $A$ y $opt(I)$ como el valor de la mejor solución para la instancia I:
\begin{definicion}[Aproximación estándar]
    La cualidad de un algoritmo de aproximación A viene dada por el cociente $$\rho_{A}(I)=\dfrac{m_{A}(I, S)}{opt(I)}$$
    Llamamos a $\rho_{A}$ \textbf{ratio de aproximación estándar}.
\end{definicion}
\begin{observacion}
    Nótese que para los problemas de minimización, el ratio de aproximación está en $\left[1, \infty\right)$, mientras que en los problemas de maximización se encuentra en $\left(0, 1\right]$.
\end{observacion}
\begin{definicion}[Aproximación diferencial]
La cualidad de un algoritmo de aproximación A viene dado por el cociente $$\delta_{A}(I)=\dfrac{|\omega(I) - m_{A}(I, S)|}{|\omega(I) - \beta(I)|}$$ 
Llamamos a $\delta_{A}$ \textbf{ratio de aproximación diferencial}.
\end{definicion}
\begin{observacion}
    El valor del ratio diferencial siempre se encuentra en $\left[0,1\right]$, independientemente de que se trate de un problema de minimización o de maximización.
\end{observacion}

Existen un gran número de clases de aproximabilidad comprendidas entre $\mathcal{PO}$ y $\mathcal{NPO}$. Además, no existe una relación directa entre las clases de complejidad y las clases de aproximación. Tanto es así que para cada clase de las que hablaremos existe un problema $\mathcal{NP}$-duro que pertenece a esa clase.
\begin{definicion}
    Algunas de las clases de aproximabilidad más conocidas son:
    \begin{itemize}
        \item[$\bullet$] \textbf{Exp-$\mathcal{APX}$} y \textbf{Exp-$\mathcal{DAPX}$}. Clases de problemas para los que el mejor ratio de aproximación conocido es exponencial respecto al tamaño de la instancia.
        \item[$\bullet$] 
        \textbf{Poly-$\mathcal{APX}$} y \textbf{Poly-$\mathcal{DAPX}$}. Clases de problemas para los que el mejor ratio de aproximación conocido es polinómico respecto al tamaño de la instancia.
            \item[$\bullet$] \textbf{Log-$\mathcal{APX}$} y \textbf{Log-$\mathcal{DAPX}$}. Clases de problemas para los que el mejor ratio de aproximación conocido es logarítmico en alguna base respecto al tamaño de la instancia.
        \item[$\bullet$] \textbf{$\mathcal{APX}$} y \textbf{$\mathcal{DAPX}$}. A partir de aquí comienzan las clases de aproximación más optimistas. Estas clases corresponden a los problemas cuyo ratio de aproximación es alguna constante fija.
        \item[$\bullet$] $\mathcal{PTAS}$ y $\mathcal{DPTAS}$. Clases de problemas que admiten un \textbf{esquema de aproximación de tiempo polinomial}. Un esquema de aproximación de tiempo polinomial es una secuencia de algoritmos $A_{\varepsilon}$ que consiguen un ratio de aproximación de $1 + \varepsilon$ para cada $\varepsilon > 0$ o $1 - \varepsilon$ para problemas de maximización y para la versión diferencial. Este ratio puede obtenerse en tiempo polinómico respecto al tamaño de la instancia pero puede llegar a ser exponencial respecto a $1/\varepsilon$.
        \item[$\bullet$] $\mathcal{FPTAS}$ y $\mathcal{DFPTAS}$. Clases de problemas que admiten un \textbf{esquema de aproximación de tiempo polinómico completo}. Esto significa que además de admitir un esquema de aproximación de tiempo polinómico, también se puede obtener el ratio de aproximación de $1 \pm \varepsilon$ en tiempo polinómico respecto a $1/\varepsilon$.
        \item[$\bullet$] \textbf{0-$\mathcal{DAPX}$}. Esta clase solo se comprende en el paradigma de la aproximación diferencial. Esta clase incluye los problemas para los que cualquier algoritmo computable en tiempo polinómico devuelve para alguna instancia la peor solución posible. Es decir, que su ratio de aproximación diferencial es 0 para alguna instancia.
    \end{itemize}
    Para cada clase de aproximación estándar $C$, $DC$ denota su respectiva clase de aproximación diferencial. Nótese que el ratio de aproximación se toma teniendo en cuenta el peor entre todas las instancias del problema. 
\end{definicion}

Se tiene la siguiente jerarquía entre las clases anteriores:
$$\mathcal{PO} \subseteq \mathcal{FPTAS} \subseteq \mathcal{PTAS} \subseteq \mathcal{APX} \subseteq Log-\mathcal{APX} \subseteq Poly-\mathcal{APX} \subseteq Exp-\mathcal{APX} \subseteq \mathcal{NPO}$$
$$\mathcal{PO} \subseteq \mathcal{DFPTAS} \subseteq \mathcal{DPTAS} \subseteq \mathcal{DAPX} \subseteq Log-\mathcal{DAPX} \subseteq$$ $$\subseteq Poly-\mathcal{DAPX} \subseteq Exp-\mathcal{DAPX} \subseteq \mathcal{NPO}$$
El hecho de que cualquiera de estas inclusiones sea estricta o no es también objeto de debate. De hecho, si dos de estas clases cualesquiera fuesen iguales se tendría que todas las clases son iguales y, por consiguiente, se tendría que $\mathcal{P} = \mathcal{NP}$.
\chapter{Estado de la Cuestión}
\label{cap:estadoDeLaCuestion}

En el estado de la cuestión es donde aparecen gran parte de las referencias bibliográficas del trabajo. Una de las formas más cómodas de gestionar la bibliografía en {\LaTeX} es utilizando \textbf{bibtex}. Las entradas bibliográficas deben estar en un fichero con extensión \textit{.bib} (con esta plantilla se proporciona el fichero biblio.bib, donde están las entradas referenciadas más abajo). Cada entrada bibliográfica tiene una clave que permite referenciarla desde cualquier parte del texto con los siguiente comandos:

\begin{itemize}
\item Referencia bibliografica con cite: \cite{ldesc2e}
\item Referencia bibliográfica con citep: \citep{notsoshort}
\item Referencia bibliográfica con citet: \citet{latexAPrimer}
\end{itemize}

Es posible citar más de una fuente, como por ejemplo \citep{latexCompanion,LaTeXLamport,texKnuth}

Después, \LaTeX se ocupa de rellenar la sección de bibliografía con las entradas \textbf{que hayan sido citadas} (es decir, no con todas las entradas que hay en el .bib, sino sólo con aquellas que se hayan citado en alguna parte del texto).

Bibtex es un programa separado de latex, pdflatex o cualquier otra cosa que se use para compilar los .tex, de manera que para que se rellene correctamente la sección de bibliografía es necesario compilar primero el trabajo (a veces es necesario compilarlo dos veces), compilar después con bibtex, y volver a compilar otra vez el trabajo (de nuevo, puede ser necesario compilarlo dos veces). 

\chapter{Descripción del Trabajo}
\label{cap:descripcionTrabajo}

Aquí comienza la descripción del trabajo realizado. Se deben incluir tantos capítulos como sea necesario para describir de la manera más completa posible el trabajo que se ha llevado a cabo. Como muestra la figura \ref{fig:sampleImage}, está todo por hacer.

\begin{figure}[h]
	\centering
	\includegraphics[width = 0.5\textwidth]{Imagenes/Vectorial/Todo.pdf}
	\caption{Ejemplo de imagen}
	\label{fig:sampleImage}
\end{figure}

Si te sirve de utilidad,  puedes incluir tablas para mostrar resultados, tal como se ve en la tabla \ref{tab:sampleTable}.


\begin{table}
	\centering
	\begin{tabular}{c|c|c}
		\textbf{Col 1} & \textbf{Col 2} & \textbf{Col 3} \\
		\hline\hline
		3 & 3.01 & 3.50\\
		6 & 2.12 & 4.40\\
		1 & 3.79 & 5.00\\
		2 & 4.88 & 5.30\\
		4 & 3.50 & 2.90\\
		5 & 7.40 & 4.70\\
		\hline
	\end{tabular}
	\caption{Tabla de ejemplo}
	\label{tab:sampleTable}
\end{table}

%\include{Capitulos/Capitulo4}
%\include{Capitulos/Capitulo5}
\chapter{Conclusiones y Trabajo Futuro}
\label{cap:conclusiones}

Conclusiones del trabajo y líneas de trabajo futuro.

Antes de la entrega de actas de cada convocatoria, en el plazo que se indica en el calendario de los trabajos de fin de grado, el estudiante entregará en el Campus Virtual la versión final de la memoria en PDF.




%%%%%%%%%%%%%%%%%%%%%%%%%%%%%%%%%%%%%%%%%%%%%%%%%%%%%%%%%%%%%%%%%%%%%%%%%%%
% Si el TFG se escribe en inglés, comentar las siguientes líneas 
% porque no es necesario incluir nuevamente las Conclusiones en inglés
%\begin{otherlanguage}{english}
%\chapter*{Introduction}
\label{cap:introduction}
\addcontentsline{toc}{chapter}{Introduction}

Introduction to the subject area. This chapter contains the translation of Chapter \ref{cap:introduccion}.










%\chapter*{Conclusions and Future Work}
\label{cap:conclusions}
\addcontentsline{toc}{chapter}{Conclusions and Future Work}

Conclusions and future lines of work. This chapter contains the translation of Chapter \ref{cap:conclusiones}.



%\end{otherlanguage}
%%%%%%%%%%%%%%%%%%%%%%%%%%%%%%%%%%%%%%%%%%%%%%%%%%%%%%%%%%%%%%%%%%%%%%%%%%%

\chapter*{Contribuciones Personales}
\label{cap:contribucionesPersonales}
\addcontentsline{toc}{chapter}{Contribuciones Personales}

En caso de trabajos no unipersonales, cada participante indicará en la memoria su contribución al proyecto con una extensión de al menos dos páginas por cada uno de los participantes.

En caso de trabajo unipersonal, elimina esta página en el fichero \texttt{TFGTeXiS.tex} (comenta o borra la línea \verb|\chapter*{Contribuciones Personales}
\label{cap:contribucionesPersonales}
\addcontentsline{toc}{chapter}{Contribuciones Personales}

En caso de trabajos no unipersonales, cada participante indicará en la memoria su contribución al proyecto con una extensión de al menos dos páginas por cada uno de los participantes.

En caso de trabajo unipersonal, elimina esta página en el fichero \texttt{TFGTeXiS.tex} (comenta o borra la línea \verb|\chapter*{Contribuciones Personales}
\label{cap:contribucionesPersonales}
\addcontentsline{toc}{chapter}{Contribuciones Personales}

En caso de trabajos no unipersonales, cada participante indicará en la memoria su contribución al proyecto con una extensión de al menos dos páginas por cada uno de los participantes.

En caso de trabajo unipersonal, elimina esta página en el fichero \texttt{TFGTeXiS.tex} (comenta o borra la línea \verb|\include{Capitulos/ContribucionesPersonales}|).

\section*{Estudiante 1}
Al menos dos páginas con las contribuciones del estudiante 1.

\section*{Estudiante 2}
Al menos dos páginas con las contribuciones del estudiante 2. En caso de que haya más estudiantes, copia y pega una de estas secciones.

|).

\section*{Estudiante 1}
Al menos dos páginas con las contribuciones del estudiante 1.

\section*{Estudiante 2}
Al menos dos páginas con las contribuciones del estudiante 2. En caso de que haya más estudiantes, copia y pega una de estas secciones.

|).

\section*{Estudiante 1}
Al menos dos páginas con las contribuciones del estudiante 1.

\section*{Estudiante 2}
Al menos dos páginas con las contribuciones del estudiante 2. En caso de que haya más estudiantes, copia y pega una de estas secciones.



%
% Bibliografía
%
% Si el TFM se escribe en inglés, editar TeXiS/TeXiS_bib para cambiar el
% estilo de las referencias
%---------------------------------------------------------------------
%
%                      configBibliografia.tex
%
%---------------------------------------------------------------------
%
% bibliografia.tex
% Copyright 2009 Marco Antonio Gomez-Martin, Pedro Pablo Gomez-Martin
%
% This file belongs to the TeXiS manual, a LaTeX template for writting
% Thesis and other documents. The complete last TeXiS package can
% be obtained from http://gaia.fdi.ucm.es/projects/texis/
%
% Although the TeXiS template itself is distributed under the 
% conditions of the LaTeX Project Public License
% (http://www.latex-project.org/lppl.txt), the manual content
% uses the CC-BY-SA license that stays that you are free:
%
%    - to share & to copy, distribute and transmit the work
%    - to remix and to adapt the work
%
% under the following conditions:
%
%    - Attribution: you must attribute the work in the manner
%      specified by the author or licensor (but not in any way that
%      suggests that they endorse you or your use of the work).
%    - Share Alike: if you alter, transform, or build upon this
%      work, you may distribute the resulting work only under the
%      same, similar or a compatible license.
%
% The complete license is available in
% http://creativecommons.org/licenses/by-sa/3.0/legalcode
%
%---------------------------------------------------------------------
%
% Fichero  que  configura  los  parámetros  de  la  generación  de  la
% bibliografía.  Existen dos  parámetros configurables:  los ficheros
% .bib que se utilizan y la frase célebre que aparece justo antes de la
% primera referencia.
%
%---------------------------------------------------------------------


%%%%%%%%%%%%%%%%%%%%%%%%%%%%%%%%%%%%%%%%%%%%%%%%%%%%%%%%%%%%%%%%%%%%%%
% Definición de los ficheros .bib utilizados:
% \setBibFiles{<lista ficheros sin extension, separados por comas>}
% Nota:
% Es IMPORTANTE que los ficheros estén en la misma línea que
% el comando \setBibFiles. Si se desea utilizar varias líneas,
% terminarlas con una apertura de comentario.
%%%%%%%%%%%%%%%%%%%%%%%%%%%%%%%%%%%%%%%%%%%%%%%%%%%%%%%%%%%%%%%%%%%%%%
\setBibFiles{%
biblio%
}

%%%%%%%%%%%%%%%%%%%%%%%%%%%%%%%%%%%%%%%%%%%%%%%%%%%%%%%%%%%%%%%%%%%%%%
% Definición de la frase célebre para el capítulo de la
% bibliografía. Dentro normalmente se querrá hacer uso del entorno
% \begin{FraseCelebre}, que contendrá a su vez otros dos entornos,
% un \begin{Frase} y un \begin{Fuente}.
%
% Nota:
% Si no se quiere cita, se puede eliminar su definición (en la
% macro setCitaBibliografia{} ).
%%%%%%%%%%%%%%%%%%%%%%%%%%%%%%%%%%%%%%%%%%%%%%%%%%%%%%%%%%%%%%%%%%%%%%
\setCitaBibliografia{
\begin{FraseCelebre}
\begin{Frase}
  Y así, del mucho leer y del poco dormir, se le secó el celebro de
  manera que vino a perder el juicio.\\ 
  \textcolor{red}{(modificar en Cascaras$\backslash$bibliografia.tex)}
\end{Frase}
\begin{Fuente}
  Miguel de Cervantes Saavedra
\end{Fuente}
\end{FraseCelebre}
}

%%
%% Creamos la bibliografia
%%
\makeBib

% Variable local para emacs, para  que encuentre el fichero maestro de
% compilación y funcionen mejor algunas teclas rápidas de AucTeX

%%%
%%% Local Variables:
%%% mode: latex
%%% TeX-master: "../Tesis.tex"
%%% End:



% Apéndices
\appendix
\chapter{Título del Apéndice A}
\label{Appendix:Key1}

Los apéndices son secciones al final del documento en las que se agrega texto con el objetivo de ampliar los contenidos del documento principal.
\chapter{Título del Apéndice B}
\label{Appendix:Key2}

Se pueden añadir los apéndices que se consideren oportunos.
%\include{Apendices/appendixC}
%\include{...}
%\include{...}
%\include{...}
\backmatter



%
% Índice de palabras
%

% Sólo  la   generamos  si  está   declarada  \generaindice.  Consulta
% TeXiS.sty para más información.

% En realidad, el soporte para la generación de índices de palabras
% en TeXiS no está documentada en el manual, porque no ha sido usada
% "en producción". Por tanto, el fichero que genera el índice
% *no* se incluye aquí (está comentado). Consulta la documentación
% en TeXiS_pream.tex para más información.
\ifx\generaindice\undefined
\else
%%---------------------------------------------------------------------
%
%                        TeXiS_indice.tex
%
%---------------------------------------------------------------------
%
% TeXiS_indice.tex
% Copyright 2009 Marco Antonio Gomez-Martin, Pedro Pablo Gomez-Martin
%
% This file belongs to TeXiS, a LaTeX template for writting
% Thesis and other documents. The complete last TeXiS package can
% be obtained from http://gaia.fdi.ucm.es/projects/texis/
%
% This work may be distributed and/or modified under the
% conditions of the LaTeX Project Public License, either version 1.3
% of this license or (at your option) any later version.
% The latest version of this license is in
%   http://www.latex-project.org/lppl.txt
% and version 1.3 or later is part of all distributions of LaTeX
% version 2005/12/01 or later.
%
% This work has the LPPL maintenance status `maintained'.
% 
% The Current Maintainers of this work are Marco Antonio Gomez-Martin
% and Pedro Pablo Gomez-Martin
%
%---------------------------------------------------------------------
%
% Contiene  los  comandos  para  generar  el índice  de  palabras  del
% documento.
%
%---------------------------------------------------------------------
%
% NOTA IMPORTANTE: el  soporte en TeXiS para el  índice de palabras es
% embrionario, y  de hecho  ni siquiera se  describe en el  manual. Se
% proporciona  una infraestructura  básica (sin  terminar)  para ello,
% pero  no ha  sido usada  "en producción".  De hecho,  a pesar  de la
% existencia de  este fichero, *no* se incluye  en Tesis.tex. Consulta
% la documentación en TeXiS_pream.tex para más información.
%
%---------------------------------------------------------------------


% Si se  va a generar  la tabla de  contenidos (el índice  habitual) y
% también vamos a  generar el índice de palabras  (ambas decisiones se
% toman en  función de  la definición  o no de  un par  de constantes,
% puedes consultar modo.tex para más información), entonces metemos en
% la tabla de contenidos una  entrada para marcar la página donde está
% el índice de palabras.

\ifx\generatoc\undefined
\else
   \addcontentsline{toc}{chapter}{\indexname}
\fi


% Generamos el índice
\printindex

% Variable local para emacs, para  que encuentre el fichero maestro de
% compilación y funcionen mejor algunas teclas rápidas de AucTeX

%%%
%%% Local Variables:
%%% mode: latex
%%% TeX-master: "./tesis.tex"
%%% End:

\fi

%
% Lista de acrónimos
%

% Sólo  lo  generamos  si  está declarada  \generaacronimos.  Consulta
% TeXiS.sty para más información.


\ifx\generaacronimos\undefined
\else
%---------------------------------------------------------------------
%
%                        TeXiS_acron.tex
%
%---------------------------------------------------------------------
%
% TeXiS_acron.tex
% Copyright 2009 Marco Antonio Gomez-Martin, Pedro Pablo Gomez-Martin
%
% This file belongs to TeXiS, a LaTeX template for writting
% Thesis and other documents. The complete last TeXiS package can
% be obtained from http://gaia.fdi.ucm.es/projects/texis/
%
% This work may be distributed and/or modified under the
% conditions of the LaTeX Project Public License, either version 1.3
% of this license or (at your option) any later version.
% The latest version of this license is in
%   http://www.latex-project.org/lppl.txt
% and version 1.3 or later is part of all distributions of LaTeX
% version 2005/12/01 or later.
%
% This work has the LPPL maintenance status `maintained'.
% 
% The Current Maintainers of this work are Marco Antonio Gomez-Martin
% and Pedro Pablo Gomez-Martin
%
%---------------------------------------------------------------------
%
% Contiene  los  comandos  para  generar  el listado de acrónimos
% documento.
%
%---------------------------------------------------------------------
%
% NOTA IMPORTANTE:  para que la  generación de acrónimos  funcione, al
% menos  debe  existir  un  acrónimo   en  el  documento.  Si  no,  la
% compilación  del   fichero  LaTeX  falla  con   un  error  "extraño"
% (indicando  que  quizá  falte  un \item).   Consulta  el  comentario
% referente al paquete glosstex en TeXiS_pream.tex.
%
%---------------------------------------------------------------------


% Redefinimos a español  el título de la lista  de acrónimos (Babel no
% lo hace por nosotros esta vez)

\def\listacronymname{Lista de acrónimos}

% Para el glosario:
% \def\glosarryname{Glosario}

% Si se  va a generar  la tabla de  contenidos (el índice  habitual) y
% también vamos a  generar la lista de acrónimos  (ambas decisiones se
% toman en  función de  la definición  o no de  un par  de constantes,
% puedes consultar config.tex  para más información), entonces metemos
% en la  tabla de contenidos una  entrada para marcar  la página donde
% está el índice de palabras.

\ifx\generatoc\undefined
\else
   \addcontentsline{toc}{chapter}{\listacronymname}
\fi


% Generamos la lista de acrónimos (en realidad el índice asociado a la
% lista "acr" de GlossTeX)

\printglosstex(acr)

% Variable local para emacs, para  que encuentre el fichero maestro de
% compilación y funcionen mejor algunas teclas rápidas de AucTeX

%%%
%%% Local Variables:
%%% mode: latex
%%% TeX-master: "../Tesis.tex"
%%% End:

\fi

%
% Final
%
%---------------------------------------------------------------------
%
%                      fin.tex
%
%---------------------------------------------------------------------
%
% fin.tex
% Copyright 2009 Marco Antonio Gomez-Martin, Pedro Pablo Gomez-Martin
%
% This file belongs to the TeXiS manual, a LaTeX template for writting
% Thesis and other documents. The complete last TeXiS package can
% be obtained from http://gaia.fdi.ucm.es/projects/texis/
%
% Although the TeXiS template itself is distributed under the 
% conditions of the LaTeX Project Public License
% (http://www.latex-project.org/lppl.txt), the manual content
% uses the CC-BY-SA license that stays that you are free:
%
%    - to share & to copy, distribute and transmit the work
%    - to remix and to adapt the work
%
% under the following conditions:
%
%    - Attribution: you must attribute the work in the manner
%      specified by the author or licensor (but not in any way that
%      suggests that they endorse you or your use of the work).
%    - Share Alike: if you alter, transform, or build upon this
%      work, you may distribute the resulting work only under the
%      same, similar or a compatible license.
%
% The complete license is available in
% http://creativecommons.org/licenses/by-sa/3.0/legalcode
%
%---------------------------------------------------------------------
%
% Contiene la última página
%
%---------------------------------------------------------------------


% Ponemos el marcador en el PDF
\ifpdf
   \pdfbookmark{Fin}{fin}
\fi

\thispagestyle{empty}\mbox{}

Este texto se puede encontrar en el fichero Cascaras/fin.tex. Si deseas eliminarlo, basta con comentar la línea correspondiente al final del fichero TFGTeXiS.tex.

\vspace*{4cm}

\small

\hfill \emph{--¿Qué te parece desto, Sancho? -- Dijo Don Quijote --}

\hfill \emph{Bien podrán los encantadores quitarme la ventura,}

\hfill \emph{pero el esfuerzo y el ánimo, será imposible.}

\hfill 

\hfill \emph{Segunda parte del Ingenioso Caballero} 

\hfill \emph{Don Quijote de la Mancha}

\hfill \emph{Miguel de Cervantes}

\vfill%space*{4cm}

\hfill \emph{--Buena está -- dijo Sancho --; fírmela vuestra merced.}

\hfill \emph{--No es menester firmarla -- dijo Don Quijote--,}

\hfill \emph{sino solamente poner mi rúbrica.}

\hfill 

\hfill \emph{Primera parte del Ingenioso Caballero} 

\hfill \emph{Don Quijote de la Mancha}

\hfill \emph{Miguel de Cervantes}


\newpage
\thispagestyle{empty}\mbox{}

\newpage

% Variable local para emacs, para  que encuentre el fichero maestro de
% compilación y funcionen mejor algunas teclas rápidas de AucTeX

%%%
%%% Local Variables:
%%% mode: latex
%%% TeX-master: "../Tesis.tex"
%%% End:

%\end{otherlanguage}
\end{document}
